\section{Introdução}
\label{sec:panorama-intro}

Nas últimas duas décadas, a odontologia experimentou uma transformação tecnológica e epistemológica comparável à revolução industrial do século XIX, migrando de um modelo eminentemente artesanal, analógico e operador-dependente para um ecossistema digital integrado e balizado por dados multidimensionais \cite{mdpi2025convergence}. O que historicamente dependia em caráter exclusivo da percepção tátil do clínico, de registros físicos suscetíveis à degradação material e de técnicas laboratoriais manuais, hoje encontra sua validação e execução ancoradas em sensores ópticos submicrométricos, algoritmos de reconstrução geométrica tridimensional e plataformas de planejamento virtual avançadas \cite{roy2025applications}. 

\destaque{A ontologia do paciente deixou de ser estritamente biológica para assumir um caráter híbrido, no qual a dimensão informacional (os dados) e a dimensão material (os tecidos) interagem de forma recíproca e ininterrupta.}

O molde de alginato ou silicone, utilizado por mais de um século como o padrão-ouro inquestionável para o registro reológico das arcadas dentárias, foi progressivamente substituído e mesmo superado pelo scanner intraoral. Esta tecnologia fotogramétrica ou confocal\footnote{A tecnologia confocal utiliza um diafragma com pinhole posicionado em um plano conjugado com o ponto focal da lente objetiva para eliminar luz fora de foco. Em scanners intraorais, isso permite medir a profundidade e reconstruir a topografia tridimensional da superfície dentária sem a necessidade de aplicar pó opaco refletor sobre os dentes.} captura miríades de pontos por segundo (nuvens de pontos ou \textit{point clouds}\footnote{Nuvens de pontos (\textit{point clouds}) são conjuntos tridimensionais de coordenadas cartesianas $(X, Y, Z)$ obtidos a partir de sensores ópticos ou varreduras laser. Elas representam a forma bruta inicial do objeto escaneado antes de sofrer triangulação (geralmente por algoritmos de Delaunay) para a geração de malhas poligonais (meshes) compostas por vértices, arestas e faces.}), processando-os estocasticamente para constituir malhas poligonais (meshes) com precisão espacial da ordem de dezenas de micrômetros \cite{elsaid2025digital}. De maneira correlata, a radiografia convencional, estruturalmente limitada a projeções bidimensionais com ampliação e distorção inerentes ao feixe cônico primitivo, cedeu espaço à tomografia computadorizada de feixe cônico (CBCT). Este maquinário é capaz de gerar volumes voxelados isotrópicos\footnote{Um voxel (abreviação de \textit{volume element}) isotrópico possui dimensões perfeitamente idênticas nos três eixos espaciais $(x, y, z)$. Diferente de imagens médicas tradicionais em que os voxels podem ser anisotrópicos (alongados em um eixo), a isotropia na tomografia CBCT odontológica garante que a resolução espacial e a precisão dimensional sejam constantes, independentemente do plano anatômico de corte selecionado (axial, coronal ou sagital).} com resolução submilimétrica, destituídos de superposições anatômicas indesejadas \cite{li2025impacts}.

Mais do que o mero abandono de substratos analógicos — como fotografias impressas, modelos de gesso e padrões em cera —, essa nova matriz operativa deu lugar a plataformas de alto nível heurístico. Elas amalgamam o diagnóstico por imagem estrutural, a simulação biomecânica de vetores de força e a manufatura aditiva/subtrativa dentro de um mesmo \textit{pipeline} digital ou fluxo de trabalho integrado \cite{techsoc2025precision}.

Essa revolução digital não apenas incrementou a precisão estática e a previsibilidade clínica dos procedimentos reconstrutivos e cirúrgicos, mas também — e talvez de forma mais disruptiva a longo prazo — forjou a espinha dorsal infraestrutural necessária para a emergência dos gêmeos digitais (\textit{Digital Twins}) na odontologia, que constituem o núcleo axiomático deste livro \cite{katsoulakis2024digital}. Conforme será explorado em profundidade no \autoref{ch:fundamentos}, o conceito de gêmeo digital pressupõe a existência indissociável de um modelo virtual dinâmico, preditivo e bidirecionalmente conectado ao seu correspondente físico \cite{robles2023opentwins}. Sem a quantificação e digitalização massiva do biotipo do paciente, esse arcabouço permaneceria restrito ao domínio da especulação teórica. 

O presente capítulo oferece, portanto, um panorama denso e abrangente da odontologia digital contemporânea, perscrutando sua evolução histórica por meio de inflexões tecnológicas, o estado da arte do ferramental atual, o fluxo de trabalho digital integrado em sua completude, o papel descentralizador da tele-odontologia em rede e os desafios epistêmicos e operacionais ainda remanescentes para a capilarização plena deste novo paradigma nas ciências da saúde \cite{khoshfekr2025digital}.
